% F082｜05 易错点
\clearpage
\begin{trapbox}
\textbf{\textcolor{CTrap}{易错一：把距离和型 F012 的结论直接套到距离差型}}
\begin{center}\begin{tikzpicture}[font=\small,>=stealth,line join=round]
\draw[gray!40,rounded corners=3pt] (0,0) rectangle (14.2,8.6);\node at (7.1,8.0) {易错对照：求和与求差，取等时动点的位置不同};
\begin{scope}[shift={(2.0,4.3)},x=0.85cm,y=0.85cm]\draw[->,gray] (-0.6,0)--(2.6,0) node[right]{$x$};\draw[->,gray] (0,-.35)--(0,1.8) node[left]{$y$};\draw[orange!85!black,line width=1.6pt] (0,-2)--(1.3333333333333333,0);\draw[violet,thick] (2,1)--(1.3333333333333333,0);\draw[teal!80!black,thick] (0,-2)--(2,1);\fill[teal!80!black] (0,-2) circle(2.2pt) node[left=6pt]{$A$};\fill[teal!80!black] (2,1) circle(2.2pt) node[above right=4pt]{$B$};\fill[blue] (1.3333333333333333,0) circle(2.2pt) node[below=4pt]{$P$};\end{scope}\draw[gray] (2,2.4)--(2,4.3);\begin{scope}[shift={(8.0,2.6)},x=0.85cm,y=0.85cm]\draw[->,gray] (-0.6,0)--(4.6,0) node[right]{$x$};\draw[->,gray] (0,-.35)--(0,2.8) node[left]{$y$};\draw[orange!85!black,line width=1.6pt] (0,2)--(4,0);\draw[violet,thick] (2,1)--(4,0);\draw[teal!80!black,thick] (0,2)--(2,1);\fill[teal!80!black] (0,2) circle(2.2pt) node[left=6pt]{$A$};\fill[teal!80!black] (2,1) circle(2.2pt) node[above right=4pt]{$B$};\fill[blue] (4,0) circle(2.2pt) node[below=4pt]{$P$};\end{scope}\node[align=center] at (4.0,6.8) {求和：两定点在轴两侧\\$P$ 在线段内部};\node[align=center] at (10.4,6.8) {求差：两定点在轴同侧\\$P$ 在线段外部};\node at (4,1.4){$PA+PB=AB$};\node at (10.4,1.4){$PA-PB=AB$};
\node[text width=12.8cm,align=center] at (7.1,.65) {两幅图横纵同尺度；左图 $A(0,-2),B(2,1)$，右图 $A(0,2),B(2,1)$。};
\end{tikzpicture}\end{center}



距离和要让两段路拉成一条直路，常把两个定点放在 $x$ 轴上下两侧，得到的竖直距离是 $a+b$。

距离差应让两定点位于\textbf{同一侧}，比较两段路相差多少，得到的是 $|a-b|$。
\[
 \boxed{D=\sqrt{(m-n)^2+(a-b)^2}},
 \quad\text{不能错写成 }\sqrt{(m-n)^2+(a+b)^2}.
\]

\end{trapbox}
\clearpage
\begin{trapbox}
\textbf{\textcolor{CTrap}{易错二：从 $|f(x)|\leq D$ 就直接说值域 $[-D,D]$}}
\begin{center}\begin{tikzpicture}[font=\small,>=stealth,line join=round]
\draw[gray!40,rounded corners=3pt] (0,0) rectangle (14.2,10.2);\node at (7.1,9.6) {第一项高度更大：先增后减，只有最大值能取到};
\node[blue] at (7.1,8.65) {$\displaystyle f(x)=\sqrt{x^2+4}-\sqrt{(x-2)^2+1}$};\begin{scope}[shift={(3.1,4.6)},x=0.36cm,y=1cm]\draw[->,gray] (-6,0)--(14.5,0) node[right]{$x$};\draw[->,gray] (0,-2.8)--(0,3.4000000000000004) node[left]{$y$};\draw[gray,dashed] (-6,-2)--(14,-2);\node[left=3pt,fill=white,inner sep=1pt] at (0,-2) {$-2$};\draw[gray,dashed] (-6,2)--(14,2);\node[left=3pt,fill=white,inner sep=1pt] at (0,2) {$2$};\draw[blue,very thick,domain=-6:14,samples=240] plot(\x,{sqrt(\x*\x+4)-sqrt((\x-2)^2+1)});\draw[gray,dotted] (4,0)--(4,2.23606797749979);\fill[blue] (4,2.23606797749979) circle(2.3pt);\node[below=4pt] at (4,0){$4$};\node[blue,above=9pt] at (4,2.23606797749979) {$(4,\sqrt5)$};\end{scope}\node[anchor=west,align=left] at (9,4.5) {$a=2>b=1$\\[12pt]左端：趋近 $-2$\\[10pt]右端：趋近 $2$\\[12pt]最大值：$\sqrt5$\\[12pt]完整值域：$(-2,\sqrt5]$};
\node[text width=12.8cm,align=center] at (7.1,.65) {虚线是无限远处靠近的高度；不是有限定义域的端点。};
\end{tikzpicture}\end{center}



这仅仅说明函数不会跑到范围外，\textbf{不代表两个端点都能取到，也不代表里面所有数都能取到}。例如
\[
 \sqrt{x^2+4}-\sqrt{(x-2)^2+1}
 \in(-2,\sqrt5],
\]
而不是 $[-\sqrt5,\sqrt5]$。

\textbf{\textcolor{CTrap}{易错三：忘记减法有先后顺序}}

$PA-PB$ 和 $PB-PA$ 互为相反数；把两个根号换顺序，最大值和最小值也会对调。不能只凭 $AB$ 的长度来断定哪个是正、哪个是负。

\textbf{\textcolor{CTrap}{易错四：三点共线就一定能取到 $AB$？}}

距离差等于 $\pm AB$ 时，$P$ 要在\textbf{线段 $AB$ 的外侧}，而不是在线段中间。若 $P$ 在中间，则有 $PA+PB=AB$，对应的是\textbf{距离和}，不是差的极值。
\end{trapbox}
\clearpage
\begin{trapbox}
\textbf{\textcolor{CTrap}{易错五：高度相同还套 $x_0$ 公式}}
\begin{center}\begin{tikzpicture}[font=\small,>=stealth,line join=round]
\draw[gray!40,rounded corners=3pt] (0,0) rectangle (14.2,9.2);\node at (7.1,8.6) {高度相同：一直增大，两端都只能靠近};
\node[blue] at (7.1,7.65) {$\displaystyle h(x)=\sqrt{x^2+4}-\sqrt{(x-3)^2+4}$};\begin{scope}[shift={(4.1,4.0)},x=0.3130434782608696cm,y=0.75cm]\draw[->,gray] (-10,0)--(13.5,0) node[right]{$x$};\draw[->,gray] (0,-3.6)--(0,3.8000000000000003) node[left]{$y$};\draw[gray,dashed] (-10,-3)--(13,-3);\node[left=3pt,fill=white,inner sep=1pt] at (0,-3) {$-3$};\draw[gray,dashed] (-10,3)--(13,3);\node[left=3pt,fill=white,inner sep=1pt] at (0,3) {$3$};\draw[blue,very thick,domain=-10:13,samples=240] plot(\x,{sqrt(\x*\x+4)-sqrt((\x-3)^2+4)});\end{scope}\node[anchor=west,align=left] at (9,4.5) {$a=b=2,\ m=0,n=3$\\[12pt]对每个实数 $x$\\[10pt]都有 $-3<h(x)<3$\\[12pt]$h(3/2)=0$\\[12pt]完整值域：$(-3,3)$};
\node[text width=12.8cm,align=center] at (7.1,.65) {两条虚线均不属于函数图像；同高度且不同中心时没有最大、最小值。};
\end{tikzpicture}\end{center}



\[
 x_0=\frac{bm-an}{b-a}
\]
\textbf{只在 $a\ne b$ 时才能使用}。当 $a=b$，分母为 $0$，不是“交点在无穷远处也能代入”，而是需要重新判断：
\begin{itemize}[leftmargin=*,itemsep=.25em]
\item $m\ne n$，值域 $(-|m-n|,|m-n|)$；两端都取不到。
\item $m=n$，两根号完全一样，值域 $\{0\}$。
\end{itemize}

\textbf{\textcolor{CTrap}{易错六：把“能无限靠近”误认为“能取到”}}

例如 $x\to+\infty$ 时 $f(x)\to2$，只说明函数值越来越接近 $2$，\textbf{不等于存在一个真实的 $x$ 使 $f(x)=2$}。验证是否取到，必须看等号所要求的点能否找到，或者直接求解对应方程。

\end{trapbox}
\clearpage
\begin{trapbox}
\textbf{\textcolor{CTrap}{易错七：题目限制 $x$，却照抄全实数的值域}}
\begin{center}\begin{tikzpicture}[font=\small,>=stealth,line join=round]
\draw[gray!40,rounded corners=3pt] (0,0) rectangle (14.2,8.5);\node at (7.1,7.9) {例6：只差是否允许 x=4，值域右端开闭随之改变};
\begin{scope}[shift={(-0.2,2.3)},x=.48cm,y=17cm]\draw[->,gray] (3.4,0)--(12,0) node[right]{$x$};\draw[->,gray] (4,-.015)--(4,.28);\draw[gray,dashed] (3.5,0)--(11.5,0);\draw[blue,thick,domain=4.001:11.5,samples=130] plot(\x,{sqrt(\x*\x+4)-sqrt((\x-2)^2+1)-2});\draw[blue,fill=white] (4,{sqrt(5)-2}) circle(2.5pt);\node[below=4pt] at (4,0){$4$};\node[left=4pt] at (4,0){$2$};\node[blue,above right=8pt] at (4,{sqrt(5)-2}) {$\sqrt5$};\end{scope}\node at (3.1999999999999997,1.35) {$x>4$：$(2,\sqrt5)$};\begin{scope}[shift={(6.8,2.3)},x=.48cm,y=17cm]\draw[->,gray] (3.4,0)--(12,0) node[right]{$x$};\draw[->,gray] (4,-.015)--(4,.28);\draw[gray,dashed] (3.5,0)--(11.5,0);\draw[blue,thick,domain=4.001:11.5,samples=130] plot(\x,{sqrt(\x*\x+4)-sqrt((\x-2)^2+1)-2});\draw[blue,fill=blue] (4,{sqrt(5)-2}) circle(2.5pt);\node[below=4pt] at (4,0){$4$};\node[left=4pt] at (4,0){$2$};\node[blue,above right=8pt] at (4,{sqrt(5)-2}) {$\sqrt5$};\end{scope}\node at (10.2,1.35) {$x\ge4$：$(2,\sqrt5]$};
\node[text width=12.8cm,align=center] at (7.1,.65) {两图同尺度，纵轴从 $2$ 起显示实际函数值；下方虚线 $y=2$ 永远取不到。};
\end{tikzpicture}\end{center}



只有 $x\in\mathbb R$ 才可直接使用 F082 总表。若题目给 $x>c$、$x\in[L,R]$，先找 $x_0$，判断它在不在范围中；然后结合函数在左右哪边增大、哪边减小，分别检查端点开闭。

\textbf{\textcolor{CTrap}{易错八：根号内常数不是正平方，也硬套距离法}}

只有能写成 $(x-m)^2+a^2$ 且 $a>0$ 的形式，才能使用本条完整结论。若配方得到 $(x-m)^2-4$，不能把 $-4$ 看作固定竖直距离的平方。若 $a$ 或 $b$ 恰好是 $0$，会出现绝对值，等号位置也可能有多种情况，必须另外处理。
\end{trapbox}
\clearpage
\begin{trapbox}
\textbf{\textcolor{CTrap}{再记一个小心之处：翻折不改变某个根号，却可能改变两定点的距离}}
\begin{center}\begin{tikzpicture}[font=\small,>=stealth,line join=round]
\draw[gray!40,rounded corners=3pt] (0,0) rectangle (14.2,9.2);\node at (7.1,8.6) {翻折不改变某一根号，却会改变两定点之间的距离};
\begin{scope}[shift={(1.9,4.1)},x=1.0cm,y=1.0cm]\draw[->,gray] (-0.6,0)--(4.6,0) node[right]{$x$};\draw[->,gray] (0,-.35)--(0,2.8) node[left]{$y$};\draw[orange!85!black,line width=1.6pt] (0,2)--(4,0);\draw[violet,thick] (2,1)--(4,0);\draw[teal!80!black,thick] (0,2)--(2,1);\fill[teal!80!black] (0,2) circle(2.2pt) node[above left=4pt]{$A(0,2)$};\fill[teal!80!black] (2,1) circle(2.2pt) node[above right=4pt]{$B(2,1)$};\fill[blue] (4,0) circle(2.2pt) node[below=4pt]{$P(4,0)$};\end{scope}\begin{scope}[shift={(1.9,4.1)}]\draw[gray,dashed] (0,-2)--(4,0);\draw[gray,dotted] (0,-2)--(0,2);\draw[gray,thick,dashed] (0,-2)--(2,1);\fill[gray] (0,-2) circle(2pt) node[below right=3pt]{$A^{\prime}(0,-2)$};\end{scope}\node[anchor=west,align=left] at (8.2,4.5) {$PA^{\prime}=PA$\\[12pt]但 $AB=\sqrt5$\\[10pt]而 $A^{\prime}B=\sqrt{13}$\\[12pt]求差用同侧点的 $AB$。};
\node[text width=12.8cm,align=center] at (7.1,.65) {橙、灰两段到轴上同一点的距离相等；同侧和异侧的定点距离不能混用。};
\end{tikzpicture}\end{center}



当 $P$ 在 $x$ 轴上，$A(m,a)$ 与 $A'(m,-a)$ 到 $P$ 的距离相同，所以原函数保持不变。

但是 $AB$ 与 $A'B$ 不是同一段长度：
\[
 AB=\sqrt{(m-n)^2+(a-b)^2},\qquad
 A'B=\sqrt{(m-n)^2+(a+b)^2}.
\]
求差时用\textbf{同侧两定点}的距离 $AB$。不能把 F012 画出的异侧两点距离 $A'B$ 当成 F082 的最大距离差。

\textbf{一句防错口令：}先分高度，后定正负；检查共线，最后判开闭。
\end{trapbox}
